\section{Introduction}
The reliability of a large language model (LLM) service depends not only on the model, but also on the inference engine beneath it. Every request passes through a long execution path that covers model preparation, runtime orchestration, and hardware execution. Modern engines make this path faster through advanced scheduling, memory management, heterogeneous execution, and distributed serving~\cite{yu2022orca,kwon2023pagedattention,aminabadi2022deepspeed,zheng2024sglang,zhong2024distserve,park2026inferenceengines}. They must also support a growing range of models, optimization methods, and deployment environments. As a result, an inference engine is no longer a simple model runtime. It is a large systems stack whose layers must evolve together while preserving correctness, availability, and performance across environments.

This combination of optimization and broad compatibility makes inference-engine bugs both harmful and difficult to maintain. A faulty engine may silently corrupt an output, disrupt a service through a runtime or resource failure, or break only in a specific deployment environment. These failures are also difficult to diagnose. An incorrect output can easily be mistaken for a weakness of the model, while a latency increase may be blamed on the workload or hardware. Even after a failure is confirmed as an engine bug, its visible symptom may be far from its source. The fault may reside in the computation layer, the runtime control plane, or a deployment interface. Its repair may require either a one-line correction or a coordinated change across several modules. A symptom taxonomy can describe how an engine fails, but reliability is improved through decisions about code: where to test, where to search for a fault, and how broadly to validate a fix. Making these decisions requires connecting observed failures to the code regions that cause them and to the changes needed to repair them. Such connections turn bug observations into actionable guidance for testing and maintenance.

Prior work has studied bugs in deep learning applications, frameworks, and compilers, producing taxonomies of symptoms, causes, and affected components~\cite{islam2019bugs,humbatova2020taxonomy,yang2022frameworkbugs,chen2023frameworkbugs,shen2021compilerbugs,du2021compilerbugs}. General software-engineering studies have also shown that fixes follow recurring patterns and that faults often concentrate in a small part of a codebase~\cite{pan2009bugfixpatterns,sobreira2018dissection,liu2020vulnerabilitydistribution}. Most closely related, Liu et al.~\cite{liu2026firstlook} studied 929 bugs from five LLM inference engines and characterized their failure types, repair behavior, and evolution. These studies explain \emph{what kinds of bugs occur}. However, they leave a broader maintenance question open: \emph{how is bug maintenance structured inside the codebase, and how does it evolve over time?}

Answering this question is challenging for three reasons. First, the relation between symptoms and faulty code is many-to-many. The same symptom may originate in different layers, while one layer may produce very different failures. Issue categories alone therefore cannot identify where developers should inspect or test. Second, patch size does not reveal repair intent. Similar-sized patches may perform different operations, and one fix may combine several operations across code entities. Measuring changed lines alone therefore misses both the meaning and the spread of maintenance work. Third, inference engines evolve rapidly and at different rates. Changes in project size or bug mix can obscure whether repairs themselves are reaching a wider code surface or shifting their architectural focus. Addressing these challenges requires a joint analysis of failure symptoms, code locations, and repair patches across projects and over time.

In this paper, we take a code-centered view of maintenance in LLM inference engines. We connect each observed failure to two forms of code evidence: where the fault resides and how developers repair it. The first captures the affected architectural layer and code entities. The second captures the repair operation and repair scope, including patch breadth and architectural distribution. Rather than report these dimensions in isolation, we follow each bug from its external effect to its internal location and then to the resulting code change. By studying these connections across projects and over time, we reveal where maintenance effort accumulates, what developers repeatedly change, and how repair boundaries evolve.

We use this view to answer three research questions:
\begin{itemize}
    \item \textbf{RQ1---Failure localization:} Across LLM inference engines, how are observed failure symptoms associated with the architectural layers in which faults reside?
    \item \textbf{RQ2---Repair structure:} How is bug-fixing work structured across code entities and architectural layers?
    \item \textbf{RQ3---Repair evolution:} How does the scope of bug fixes evolve over time?
\end{itemize}

To answer these questions, we conduct the largest code-centered empirical study of LLM inference-engine bugs to date. We study six widely used open-source engines: llama.cpp, vLLM, DeepSpeed, MLC-LLM, TensorRT-LLM, and Text Generation Inference (TGI). Together, they cover cloud and local deployment, single and multi-device execution, and a wide range of hardware and software backends. We collect all 50,121 bug-fix-related pull requests submitted to their official repositories from the start of each project to March 2026. We then select the pull requests merged into the projects' codebases and manually review each one to determine whether it implements a bug fix. This process yields a study corpus of 12,707 merged and manually verified bug-fixing pull requests. For each fix, we combine repository metadata, development discussions, and the complete code patch. We classify failures into ten symptom types, locate faulty code in nine architectural layers, and identify nine types of repair operations. We then map the fixes to code entities, measure repair scope through patch breadth and architectural distribution, and track its monthly evolution. This process gives us an end-to-end view of each repair, from the failure observed by a user to the code modified by a developer.

Our study reveals three main properties of inference-engine maintenance. First, failures are strongly shaped by system architecture. Bugs that directly threaten correctness or availability, i.e., incorrect outputs, runtime crashes, and model-loading failures, account for 4,001 of the 12,707 studied fixes (31.5\%). Some layers have narrow and recognizable failure signatures. Bugs in build and release tooling mainly prevent compilation or installation, while bugs in user-facing configuration and entry points often appear as API or argument-parsing errors. By contrast, bugs in request scheduling and runtime control can affect correctness, availability, or resource use. Symptoms can therefore provide strong localization clues for specialized layers, whereas failures in the central execution logic require broader diagnosis and testing.

Second, maintenance risk is highly concentrated. The top 20\% of files account for 79.6\% of bug-fixing occurrences, and the top 20\% of functions account for 68.4\%. The pattern appears in all six engines. Hotspot files therefore contain multiple functions that repeatedly participate in bug fixes, rather than only one isolated faulty function. Repair work is concentrated as well: Logic Correction and Adaptation account for 14,702 of 19,965 repair-operation instances (73.6\%). Together, these results show that the main maintenance burden comes from a limited code surface and the continuing need to preserve correct behavior across evolving hardware and software.

Third, bug fixes are becoming less local. Most fixes are still compact. The median patch changes only 14 lines, but large repairs are not rare: roughly 10\% changes at least 165 lines. More importantly, from 2023 to early 2026, repairs increasingly spanned several files instead of remaining within one, while Hardware Adaptation occupied a growing share of bug-fixing activity. This shift means that inference-engine maintenance is moving from isolated edits toward coordinated changes across related parts of the codebase and a more prominent hardware-adaptation focus. Review and regression testing must therefore follow the full change, because checking only the first failing file can overlook related changes.

These findings change the view of inference-engine reliability in several ways. Bugs are not spread evenly across a repository; they form stable code hotspots. Symptoms are not independent of architecture; they carry different amounts of information about fault location. Repairs are not dominated by generic cleanup; they mainly correct logic and adapt behavior across environments. Finally, maintenance does not remain equally local as engines mature; the code boundary of a fix tends to widen and its architectural focus can shift. This knowledge can guide hotspot-based testing, architecture-aware debugging, and code-aware automation for LLM systems.

In summary, this paper makes the following contributions:
\begin{itemize}
    \item We present the largest code-centered empirical study of LLM inference-engine bugs, covering 12,707 merged and manually verified bug-fixing pull requests from six open-source engines. We connect failure symptoms, faulty code, and repair patches in an analysis framework.
    \item We characterize bug-fixing work in LLM inference engines. We show how system architecture shapes failure symptoms, where maintenance effort concentrates, and how repair operations differ across architectural layers.
    \item We provide a large-scale view of how inference-engine repair scope changes over time, including patch breadth and architectural distribution. Our results expose an expanding repair boundary and growing Hardware Adaptation focus, guiding testing, review, and automated repair for LLM serving systems.
\end{itemize}

\begin{tcolorbox}[
    colback=white,
    colframe=black,
    boxrule=0.5pt,
    arc=1pt,
    left=4pt,
    right=4pt,
    top=3pt,
    bottom=3pt,
    before skip=4pt,
    after skip=0pt
]
\textbf{Artifact:} Code and data are available at \url{https://github.com/Ellliot-ECNU/llm-engine-bugs}.
\end{tcolorbox}
